\documentclass[paper=a4]{scrartcl}
\usepackage[utf8]{inputenc}
\usepackage[english]{babel}
\usepackage{enumitem}
\usepackage{hyperref}

\title{IDP Project Description} 
\author{Henry He}

\begin{document}
\maketitle

\section{Title}
\subsection*{English}
SBM-based Detection of Malware Spread in a Campus Network using Central
Flow-Monitoring Data

\subsection*{Deutsch}
SBM-basierte Erkennung von Malwareausbreitung in einem Campusncetzwerk mithilfe
von Central Flow-Monitoring Daten

\section{Content}
Communication systems are considered one of the defining technologies of today's
interconnected world. They provide a platform for individuals to connect,
collaborate and share information, regardless of geographical distance.
Companies need communication networks to compete in today's fast paced business
environment as it enables higher revenues via streamlined operations, higher
efficiency and agile supply chains. 

The importance of communication networks in both industry and society has made
them a target for various attacks and breaches. These attacks can range from
hacking, malware, phishing, and cyber espionage, which can result in sensitive
information being compromised, systems being shut down, or personal privacy
being violated through unauthorized access. To combat these threats and ensure a
safe operation of communication networks, it is necessary to detect and address
attacks on networks and the sensitive information they carry as early as
possible. Effective counter strategies require a deep understanding of attack
procedures and the design of robust defense mechanisms. 

In this context, the proposed IDP focuses on the collection and analysis of
network traffic. The gained insights should then be utilized to detect malicious
behavior. As a first step, an environment shall be implemented where known
attacks can be carried out in a secure and observable manner. As a second step,
methods for identifying traffic patterns of different attacks shall be
evaluated. Here, we want to study malware spread in a small part of the campus
network using central monitoring. The acquired traffic data are aggregated in a
\textit{traffic dispersion graph} (TDG) which will be used to train so-called
\textbf{Stochastic Block Models} (SBMs) which can detect and explain specific
forms of attacks.

\section{Relation to Application Area Electrical and Computer Engineering}
Resilient network design and planning are one of the main research topics at the
Chair of Communication Systems (LKN) in the Department of Electrical and
Computer Engineering. As the volume and complexity of network traffic continues
to grow, characterizing network behavior becomes more and more difficult. This
presents a challenge for network management, which heavily relies on manual
identification of network dynamics and decision-making. The goal of this project
is to evaluate an anomaly detector using ML models, a popular tool in computer
science research. The finding of this IDP might help to overcome these
limitations and lead towards networks capable of understanding themselves and
react to threats autonomously, in turn unburden network operators.

\section{Choice of Lecture}
As communication networks form the integral part of this IDP, the lecture
``Communication Networks'' given by Prof. Kellerer teaches their foundations and
functioning from an Electrical and Computer Engineering point of view. In
particular will the proposed model analyze the network as a graph which is also
used in the lecture to assess routing algorithms as well as network planing and
management.  

\section{Milestones}
\begin{enumerate}
    \item Research on common attack patterns and detection/mitigation procedures
    \item Set up test environment for controlled attacks and design attack
    scenarios
    \item Design and implement architecture for collecting ground-truth data
    during attacks
    \item Evaluate how SBMs help detecting attack patterns in an explorative
    data analysis
\end{enumerate}

\section{Timeline}
The first 2 weeks of the project are scheduled for research described in M1.
This will be followed by the main part of the IDP, the planning and
implementation of the attack environment and data collection architecture (M2 \&
3). This part is expected to take 12-18 weeks. 2-4 weeks will be assigned to the
data analysis and SBM evaluation (M4). The remaining time of the project will be
spent for evaluation, documentation, and preparation of the presentation.

\end{document}
